\section{Inward comparison and the conclusion}\label{sec:inward}

We use the center supplied by \Cref{ce:center,ce:weights} at a point
with the density and tangent properties specified there.  If the varifold
lies on the center in a neighborhood of the point, \Cref{ce:constancy}
already gives regularity.  We therefore assume the positive-height
alternative of \Cref{ce:weights} and derive a contradiction.  In particular,
none of the arguments below assumes height doubling at the outset.
The inward subballs in \Cref{ce:weights} controlled the polynomial
weights at each fixed radius.  We now use those bounds in a separate
signed-mass comparison near the same boundary: a radial test handles a
thin outer collar, and the signed-excess theorem applies on interior balls.

\subsection{Signed comparison and frequency identities}

The comparison of height with energy follows Almgren's frequency method
\cite[Section~3.4]{DLS2011}; cutoff frequency on a center manifold is
developed in \cite[Section~3]{DLS2016III}. We derive the identities needed
here from stationarity. The signed-excess estimate supplies the additional
inequality that will yield the inward height comparison.

Keep the physical normal displacement $z$ and the radial functions
\[
 \phi_r=(1-\rho^2/r^2)_+^p,
 \qquad \psi_r=2p(1-\rho^2/r^2)_+^{p-1},
\]
where $p\geq4$ is the fixed integer chosen in \Cref{ce:weights}.
For clarity, write
\[
 \int^{(r)} f\,\dd\nu=r^{-m}\int f\,\dd\nu.
\]
This convention applies both to measures on $S$ and to measures on the
varifold; base functions in a varifold integral are pulled back by nearest
projection.  All integrals below have this normalization unless explicitly
written otherwise.  Abbreviate $\phi_r,\psi_r$ to $\phi,\psi$ at the radius
under consideration.  With $\mathbf u=\nabla_S\rho$, set
\[
 J_z=I-A_z,\qquad K_z=J_z^{-1},\qquad v=K_z\mathbf u,
 \qquad b=B[v,v],\qquad h_r=1-\frac{b}{|v|^2}.
\]
Here $A_z$ is the shape operator in the normal direction $z$; it is
distinct from the scalar $A$ defined next.  Recall that $0\leq h_r\leq1$
and that $h_r$ has no dependence on the radius at a fixed point and tangent
plane.  On the central fiber its value is immaterial, since $z=0$ there.
Define
\begin{equation}\label{in:variables}
 \begin{aligned}
 H&=\int^{(r)}\psi |z|^2h_r\,\dd\mu,
 &A&=r^2\int^{(r)}\phi\,\dd M,\\
 E&=r^2\int^{(r)}\phi\,\dd D,
 &R&=\int^{(r)}\rho^2\psi b\,\dd\mu,\\
 P&=A/H,& n_*&=E/H,\qquad q_*=R/H.
 \end{aligned}
\end{equation}
The scalar ratios $n_*,q_*$ are distinct from the ambient dimension
parameters.  The comparison we need is $2A\leq E$ up to small errors:
after division by $H$, its coefficient one gives $n_*\geq2P$ up to
those errors.  To see its role, omit the errors from the frequency
inequalities proved below.  With a dot denoting differentiation in
$\log r$, their leading terms give, for $P\geq0$,
\[
 q_*\geq n_*^2,\qquad
 \dot P\geq2P-n_*+q_*-2Pn_*
       \geq(n_*-1)(n_*-2P).
\]
Thus the signed comparison makes this leading derivative nonnegative
when $P$ is large.  We first preserve its coefficient through the
boundary of the ball; the bootstrap will then estimate the actual
errors in this calculation.

Put
\[
 \delta=r^{-1}\sup_{\rho<r}|z|,
 \qquad T=\log(r^2/H),\qquad g=T^{-20}.
\]
By \Cref{ce:weights}, after reducing the radius,
\begin{equation}\label{in:smallness}
 0<\delta\leq Cr,\qquad
 c\delta^{J_0}\leq H/r^2\leq C\delta^2,
 \qquad c|\log\delta|\leq T\leq C|\log\delta|.
\end{equation}
The finite exponent $J_0$ and all constants may depend on the fixed center.
We distinguish the unit error $\eta_1(r)=C_1\delta(r)^\alpha$ in the
collar estimate from the common error $\eta(r)=C_\eta\delta(r)^\alpha$
in the frequency argument.  Here $\alpha>0$ is fixed independently of
the multiplier and chosen sufficiently small for the finitely many power errors below, and $C_1$
depends only on the fixed center data.  We choose $C_\eta$ to dominate
the remaining error coefficients and the fixed multiples of $\eta_1$
arising from the actual collar applications.  Their multiplier norms
are verified below to have one bound independent of $r$; no factor
unbounded as $r\downarrow0$ is absorbed into $C_\eta$.  Decreasing
$\alpha$ and increasing the constants a finite number of times gives
common bounds for all the inequalities below.  In particular,
with $s=-\log r$,
\begin{equation}\label{in:exponentialerror}
 T\geq2s-C,\qquad \eta\leq C\exp(-cT).
\end{equation}

The weights first give the following comparisons:
\begin{equation}\label{in:weightcomparison}
 \begin{split}
 \int^{(r)}\phi |z|^2\,\dd\mu
   &\leq C(H+\delta^2E),\\
 \int^{(r)}\psi |z|^2\,\dd\mu
   &\leq C\{H+\delta^2(E+R)\},\\
 r^2\int^{(r)}\phi\,\dd\Lambda
   &\leq Cr^2(H+\delta^2E).
 \end{split}
\end{equation}
Indeed, $\phi\leq C\psi$, $1-h_r\leq Ce$, and $|z|\leq\delta r$
prove the first inequality.  On $\rho<r/2$ the two weights are comparable,
so the same argument applies to the second.  On $r/2\leq\rho<r$ use
$1-h_r\leq Cb$ and $r^2\psi b\leq C\rho^2\psi b$ instead.
The last inequality follows from the weighted curvature estimate in
\Cref{ce:weights} and the definition of $\Lambda$.

\begin{lemma}[The signed estimate up to the outer collar]\label{in:collar}
For $a\in C^1(B_r)$ put
\[
 \|a\|_{1,r}=\sup_{B_r}\bigl(|a|+r|\nabla_Sa|\bigr).
\]
There is a radius $r_{\mathrm c}>0$, depending only on the fixed center
data, such that for every $0<r<r_{\mathrm c}$ the first inequality below
holds for every $a\in C^1(B_r)$ with finite norm, and the second holds
independently of $a$:
\begin{equation}\label{in:collarestimates}
 \begin{split}
 r^2\left|\int^{(r)}a\phi\,\dd M\right|
   &\leq \|a\|_{1,r}\bigl[C(E+r^2H)+\eta_1 R\bigr],\\
 2A&\leq E+Cg(E+H)+Cr^2H+\eta_1 R.
 \end{split}
\end{equation}
The constants $C,C_1$ and the radius $r_{\mathrm c}$ are independent of
the multiplier.  For a family $a_r$, the coefficient of $R$ in the first
line is $\|a_r\|_{1,r}\eta_1(r)$, so each use of its smallness must
control this product.  The second line uses $a=1$ only in its thin-part
proof.
The coefficient of $E$ in the second line is exactly one before the
displayed small errors are added.
\end{lemma}

\begin{proof}
For the absolute estimate it suffices to treat $\|a\|_{1,r}\leq1$
uniformly: apply that case to $a/\|a\|_{1,r}$ when the norm is positive,
and the zero case is immediate.  The signed estimate will use the member
$a=1$ of this unit class.
Choose a small fixed $\beta>0$, with $\beta<1/4$, and put
$\sigma=\delta^\beta$; we may take $0<\alpha\leq2\beta$ in $\eta_1$.
Choose a radial cutoff $\chi$ equal to one
when $0\leq1-\rho/r\leq\sigma^2$, equal to zero when
$1-\rho/r\geq2\sigma^2$, and with derivative bounded by
$C/(\sigma^2r)$.  Write $\phi_{\rm thin}=\chi\phi$ and
$\phi_{\rm bulk}=(1-\chi)\phi$.

We first treat the bulk, including the overlap with the thin part.
Fix first an enlargement factor $c_{\mathrm{out}}>8$ sufficient for
\Cref{lt:scalar,sx:signed} and for a Caccioppoli cutoff equal to one on
$B_{8\ell}$ and supported in $B_{c_{\mathrm{out}}\ell}$.  Then choose
the fixed $c_0$ small and put $\ell=c_0\sigma^2r$.  Cover the bulk by
balls of radius $\ell$ with centers $x_i$ in its support.  These centers
have boundary depth at least $\sigma^2r$, so $c_0$ can be chosen so that
all $B_{c_{\mathrm{out}}\ell}(x_i)$ stay inside $B_r$.  Choose the cover
with bounded overlap of these enlargements and so that the intermediate
balls on which \Cref{lt:scalar} selects integers have overlaps of fixed
positive relative volume along the connecting chains.  Choose a smooth
nonnegative partition of unity $\{\vartheta_i\}$ subordinate to the
original balls, with
\[
 |D^j\vartheta_i|\leq C_j\ell^{-j}\qquad(0\leq j\leq k),
\]
where $D$ denotes base-coordinate derivatives and $k=k(m)$ is the finite
test order in \Cref{lt:potential}.  Put $\phi_i=\phi(x_i)$.  On each
enlargement $\phi$ is comparable to $\phi_i$.  The partitioned
weights, divided by $\phi_i$ and by a fixed constant,
have Lipschitz norm at most one in units of $\ell$.

The mass on each enlarged ball is at most $C\ell^m$.  To see this
without assuming a local graphical representation, observe that
$\delta r=o(\ell)$; thus the part of the tube over the ball lies in an
ambient ball of radius $C\ell$ centered at a support point.  The upper
bound follows by monotonicity, comparing with an ambient ball of radius
comparable to $r$ and using the fixed normalized mass bound at the
origin.  Empty balls need no estimate.

For the local tilt estimate, use ordinary, unnormalized physical
integrals.  Let $\xi$ be a base cutoff equal to one on $B_{4\ell}$,
supported in $B_{c_{\mathrm{out}}\ell}$, and satisfying
$|\nabla_S\xi|\leq C/\ell$.  The analogous cutoff equal to one on
$B_{8\ell}$ has the same derivative bound and outer support.
Testing stationarity with $\xi^2z$
and using the normal-displacement trace \eqref{ce:normaltrace} gives
\[
 \int\xi^2e\,\dd\mu
 \leq C\int\left[
   \bigl(|\nabla_S\xi|^2+|\mathrm{II}|^2\xi^2\bigr)|z|^2
          +\xi^2|H_S||z|\right]\,\dd\mu.
\]
Indeed, the derivative of the cutoff contributes at most
$C\xi|\nabla_S\xi||z|\sqrt e$; Young's inequality absorbs this term
and the small $|\mathrm{II}||z|e$ term from the trace.  The displayed
$|H_S||z|$ term is retained because the center need not be minimal.
The bounds $|z|\leq\delta r$ and
$|H_S|+|\nabla H_S|\leq C\delta r$, together with the local mass
bound, give
$\ell^{-m}D(B_{4\ell})\leq C[(\delta r/\ell)^2+(\delta r)^2]$.
The definition of $\Lambda$ gives
$\ell^{-m}\Lambda(B_{4\ell})\leq C(\delta r)^2$ separately.  Hence
\begin{equation}\label{in:localdensity}
 \ell^{-m}(D+\Lambda)(B_{4\ell})
       \leq C(\delta/\sigma^2)^2
       \leq C\delta^{2-4\beta}.
\end{equation}
The same estimate holds on $B_{8\ell}$, using its cutoff in
$B_{c_{\mathrm{out}}\ell}$.  These enlargement factors were fixed
before $c_0$.
The coordinates have the required fixed geometry bounds, and the terms
in $\Lambda$ still use physical units.

Apply \Cref{lt:scalar} first to select the local integers.  On the
intermediate balls just chosen, the scalar exceptional sets have relative
volume $O(\delta^{2-4\beta})$.  For small $r$ they cannot fill the
chosen overlaps of fixed positive relative volume, so their integers
agree.  The bulk is connected; call
the common integer $q_r$.  To identify it, rescale one fixed smooth
nonnegative test supported in $B_{1/2}$ to $B_{r/2}$ and normalize
its reference-volume integral to one.  Call it $\varphi_r$; then
$|D^j\varphi_r|\leq Cr^{-m-j}$ for $0\leq j\leq k$.  Partition this test
over the cover.  After pullback to a unit $\ell$-cell, each product
$\vartheta_i\varphi_r$ has $C^k$ norm at most $Cr^{-m}$, since
$\ell\leq r$.  The coordinate-volume factor in the scalar proof has
fixed scaled derivative bounds through order $k$, so multiplying by it
preserves this estimate.  The distributional estimate
\eqref{eq:lt-distribution-control} used in the proof of \Cref{lt:scalar}, with
\eqref{in:localdensity}, bounds each cell's error relative to $q_r$ by
$C\delta^{2-4\beta}\ell^m r^{-m}$.  Bounded overlap gives
$\sum_i\ell^m\leq Cr^m$, so its projected-mass average differs from
$q_r$ by at most $C\delta^{2-4\beta}$.  Tangent convergence and
the smooth rescaling of the fixed center make that same average
$Q+o(1)$.  Integrality therefore gives $q_r=Q$ for all sufficiently
small $r$.  Exact agreement on overlaps involves no error accumulated
along the chain of balls.  This also identifies the integer on the
balls used in the overlap and in the thicker error collar below.
Consequently \Cref{lt:scalar} bounds $|M|$ on each smaller ball by the
enlarged $D+\Lambda$ mass.  In particular, applying it with fixed
interior room and using \eqref{in:localdensity} also on the larger
ball gives
\begin{equation}\label{in:localscalar}
 \begin{split}
 |M|(B_{4\ell}(x_i))
   &\leq C(D+\Lambda)(B_{8\ell}(x_i)),\\
 \ell^{-m}\tau(B_{4\ell}(x_i))
   &\leq C\delta^{2-4\beta}.
 \end{split}
\end{equation}
All these balls lie inside $B_r$ by the choice of $c_0$.
Summation with comparable weights gives
\begin{equation}\label{in:bulkvariation}
 r^2\int^{(r)}\phi_{\rm bulk}\,\dd|M|
       \leq Cr^2\int^{(r)}\phi\,\dd(D+\Lambda)
       \leq C(E+r^2H),
\end{equation}
using \eqref{in:weightcomparison} to absorb a small multiple of $E$.
For every member of the unit multiplier class this also gives
$r^2|\int^{(r)}a\phi_{\rm bulk}\,\dd M|\leq C(E+r^2H)$.

For the signed estimate, invoke \Cref{sx:signed} afresh on each normal
chart at physical radius $\ell$.  The measures $M,D,\Lambda$ and the
base tests are geometric, so they give the required data in the unit
coordinates, with the measures divided by $\ell^m$.  The physical
condition below belongs to this new theorem instance; it does not change
the inherited geometry parameter in an off-center rescaling of a prior
instance.  On a cell use
\[
 d_i=\max\{C\ell^{-m}\tau(B_{4\ell}(x_i)),\ \delta^{2J_0+4}\},
 \qquad T_i=-\log d_i,
\]
with the fixed enlargement and constant sufficient for
\Cref{sx:signed}.  Equation \eqref{in:localscalar} shows that
$c|\log\delta|\leq T_i\leq(2J_0+4)|\log\delta|$, where the upper
bound follows from the floor in $d_i$.  Moreover
\[
 \ell+\sup|z|\leq C\delta^{2\beta}
                       \leq \exp(-\kappa T_i)
\]
for some fixed $\kappa>0$ and all small $r$; choose, for example,
$\kappa(2J_0+4)<2\beta$.  Thus both smallness hypotheses of
\Cref{sx:signed} hold even if $\delta$ is smaller than every fixed
power of $r$.

Apply that result to
$w_i=\vartheta_i\phi_{\rm bulk}/(C\phi_i)$ in the coordinates scaled
by $\ell$.  They are supported in $B_1$ in those coordinates and satisfy
$0\leq w_i\leq1$ and $\Lip w_i\leq1$.  Restoring
the physical measure factor $\ell^m$ and then the normalization
$r^{2-m}$ gives
\begin{equation}\label{in:bulksigned}
 \begin{split}
 2r^2\int^{(r)}\phi_{\rm bulk}\,\dd M
 &\leq r^2\int^{(r)}\phi_{\rm bulk}\,\dd D
       +C_Q r^2\int^{(r)}\phi_{\rm bulk}\,\dd\Lambda\\
 &\quad+C r^{2-m}\sum_i\phi_i\ell^m d_i e^{-F_Q(T_i)}.
 \end{split}
\end{equation}
The last sum includes the enlarged cells, even where
$\phi_{\rm bulk}$ vanishes in the overlap with the thin part.
On these cells use the local estimate \eqref{in:localscalar}:
$\phi$ is comparable to $\phi_i$ throughout $B_{8\ell}(x_i)$, and
these balls have bounded overlap inside $B_r$.  Therefore
\[
 \begin{split}
 \sum_i\phi_i\ell^m d_i
 &\leq C\sum_i\phi_i\tau(B_{4\ell}(x_i))
                  +C\delta^{2J_0+4}\sum_i\phi_i\ell^m\\
 &\leq C\sum_i\phi_i(D+\Lambda)(B_{8\ell}(x_i))
                  +C\delta^{2J_0+4}\sum_i\phi_i\ell^m\\
 &\leq Cr^m\int^{(r)}\phi\,\dd(D+\Lambda)
                  +Cr^m\delta^{2J_0+4}.
 \end{split}
\]
Here the last floor sum is bounded by the ordinary volume of the
bounded-overlap cover.  Since $T_i\geq c|\log\delta|$ and
$\exp[-F_Q(c|\log\delta|)]$ is smaller than any fixed negative power
of $T$, the remainder in \eqref{in:bulksigned} is at most
\[
 Cg\left(r^2\int^{(r)}\phi\,\dd(D+\Lambda)
                      +r^2\delta^{2J_0+4}\right).
\]
The floor term $r^2\delta^{2J_0+4}$ is at most
$C\delta^{J_0+4}H$ by \eqref{in:smallness}.  The weighted comparison
\eqref{in:weightcomparison} controls the remaining integral by
$C(E+r^2H)$ and bounds the $C_Q\Lambda$ term in
\eqref{in:bulksigned} by $Cr^2H+o(g)E$, since $r^2\delta^2=o(g)$.
This proves the second line
of \eqref{in:collarestimates} for the bulk, with its leading tilt term
unchanged.

It remains to estimate the signed mass on the thin part.  A direct
total-variation estimate here would unnecessarily add a full-tilt
term.  Instead solve a radial divergence equation.  In geodesic polar
coordinates write the volume element as $j(\rho,\omega)\,\dd\rho\,
\dd\omega$.  Throughout this thin calculation $\|a\|_{1,r}\leq1$,
and all constants are uniform over that class.  Initially put
\[
 F(\rho,\omega)=-j(\rho,\omega)^{-1}
       \int_\rho^r a(t,\omega)\phi_{\rm thin}(t)
                              j(t,\omega)\,\dd t,
 \qquad X=F\mathbf u.
\]
Then $\Div_S X=a\phi_{\rm thin}$.  Multiply $X$ by a radial cutoff
equal to one at depths $r-\rho\leq\sigma r$ and zero at depths
$r-\rho\geq2\sigma r$.  On this latter collar the integrated numerator
has size $O(r\sigma^{2p+2})$ relative to $j$.  The cutoff divergence
error $q_c$ consequently satisfies
\[
 \Div_S X=a\phi_{\rm thin}+q_c,
 \qquad |q_c|\leq C\sigma^{p+1}\phi,
\]
and is supported where $\sigma r\leq r-\rho\leq2\sigma r$.
The same integration, also differentiated in angular directions, gives
\begin{equation}\label{in:radialfield}
 |X|/r+
 |\nabla_SX-a\phi_{\rm thin}\mathbf u\otimes\mathbf u|
       \leq C\sigma^2\phi.
\end{equation}
For example, at depth $d r\leq2\sigma^2r$ the primitive divided by
$r$ is $O(d^{p+1})$, whereas $\phi\asymp d^p$; beyond that depth its
size is $O(\sigma^{2p+2})$.  On the cutoff collar the derivative loss
$\sigma^{-1}$ still leaves $\sigma^{p+1}\phi$.
This proves \eqref{in:radialfield}; bounded derivatives of the polar
volume factor only change its constants.  A radial spatial derivative
of the primitive uses the value of $a$ at its lower endpoint; an angular
derivative uses only its first spatial derivative, bounded by $1/r$.
The scale parameter $r$ is held fixed in this calculation.  Near
$\rho=r$, the preceding estimates give $|F|\leq Crd^{p+1}$ and
$|\nabla_S(F\mathbf u)|\leq Cd^p$.  Thus $X$ and its first derivative
vanish at the boundary, so extending $X$ by zero gives a $C^1$ compactly
supported base field using only values of $a$ in $B_r$.  The inner cutoff
removes the polar origin.  When $m=1$ the construction is made separately
on the two radial intervals and has no angular term.

Use the lifted field of \Cref{lt:curvedvariation}.  The $q_cM$ term
is estimated by the already proved bulk total-variation bound: for small
$\sigma$, its support has depth between $\sigma r$ and $2\sigma r$, where
$\chi=0$ and $\phi_{\rm bulk}=\phi$.  Its contribution is at most
$C\sigma^{p+1}(E+r^2H)$.  Put
$W=\int^{(r)}\phi|z|^2\,\dd\mu\leq C(H+\delta^2E)$.
Equation \eqref{in:radialfield} gives $|X|\leq C\sigma^2r\phi$ and
$|\nabla_SX|\leq C\phi$.  After multiplication by $r^2$, the three
error types in \Cref{lt:curvedvariation} are therefore bounded by
\[
 C\delta rE,\qquad C\sigma^2r^3W,\qquad
 C\sigma^2r^2\sqrt{EW},
\]
using the weighted curvature-square bound for the term containing
$\nabla^\perp H_S$.  Young's inequality puts these terms into
$Cr^2H+C(\delta r+\sigma^2)E$.  The stress left after its radial part
costs $C\sigma^2E$.  Since $p\geq4$, $r\leq1$, and
$0<\alpha\leq2\beta$, all these small coefficients of $E$ are at most
$\eta_1$ after fixing $C_1$.  The main stress term gives
\begin{equation}\label{in:thinmass}
 r^2\left|\int^{(r)}a\phi_{\rm thin}\,\dd M\right|
 \leq Cr^2\int^{(r)}\phi_{\rm thin}B[\mathbf u,\mathbf u]\,\dd\mu
                  +\eta_1 E+Cr^2H.
\end{equation}
Since $\rho\geq r/2$ on the thin part,
\[
 \phi_{\rm thin}\leq C\sigma^2(\rho/r)^2\psi,
 \qquad B[\mathbf u,\mathbf u]\leq2b+C|z|^2e.
\]
Thus the first term of \eqref{in:thinmass} is at most
$C\sigma^2R+C\delta^2r^2E\leq\eta_1(R+E)$, after increasing $C_1$.
Adding the unit bulk bound gives
\[
 r^2\left|\int^{(r)}a\phi\,\dd M\right|
 \leq C(E+r^2H)+\eta_1R\qquad(\|a\|_{1,r}\leq1).
\]
Homogeneity now gives the first line of \eqref{in:collarestimates} on
the same radius interval.  For the signed estimate use $a=1$, add twice
its thin-part bound to the bulk signed estimate, and use
$\int\phi_{\rm bulk}\,\dd D\leq\int\phi\,\dd D$.
Because $\eta_1=o(g)$ by \eqref{in:smallness}, its $E$ term is absorbed
into $CgE$ after reducing the fixed radius.  Enlarge $C_1$ once to cover
the fixed factors in both conclusions.  The coefficient one has
been preserved as claimed.
\end{proof}

\begin{lemma}[Frequency inequalities]\label{in:frequency}
With a dot denoting differentiation with respect to $\log r$, the
quantities in \eqref{in:variables} satisfy
\begin{align}
 n_*&\leq\sqrt{q_*}+Cr^2+\eta n_*,\label{in:normalfrequency}\\
 \left|\frac{\dot A}{H}-(2P-n_*+q_*)\right|
  &\leq C\{r\sqrt{q_*}+r^2(1+n_*)\}+\eta(n_*+q_*),
                         \label{in:horizontalfrequency}\\
 -m\leq\frac{\dot H}{H}
  &\leq2n_*+Cr+\eta(n_*+q_*).\label{in:heightfrequency}
\end{align}
\end{lemma}

\begin{proof}
Test normal stationarity with $z\phi$.  The exact normal-displacement trace formula
\eqref{ce:normaltrace} and \eqref{in:weightcomparison} give
\begin{equation}\label{in:normalflux}
 \left|E-\int^{(r)}\rho\psi
                     z\cdot P_{vh}v\,\dd\mu\right|
       \leq Cr^2H+\eta E.
\end{equation}
In particular the $H_S\cdot z$ term is bounded using the weighted
curvature-square estimate, rather than a height-independent error.
For the tangent projection $P_T$, the algebra is exact: $P_T^2=P_T$
and $(P_T)_{hh}=I-B$
imply
\[
 |P_{vh}v|^2=v\cdot(B-B^2)v
       =b-|Bv|^2
       \leq b-\frac{b^2}{|v|^2}=bh_r.
\]
Cauchy--Schwarz therefore bounds the absolute value of the flux in
\eqref{in:normalflux} by $\sqrt{RH}$.  Division by $H$ proves
\eqref{in:normalfrequency}.

For the horizontal identity put $Z_0=\rho\mathbf u$ on $S$ and let
$\overline Z_0(q)=J_zZ_0(\pi q)$ be its unweighted lift.  We first
compute the full derivative of the kernel in the test
$(\phi\circ\pi)\overline Z_0$.  Since $D\pi=K_zP_h$ and
$\nabla_S\phi=-(\rho/r^2)\psi\mathbf u$,
\[
 \nabla(\phi\circ\pi)=-\frac{\rho}{r^2}\psi K_z\mathbf u.
\]
Both this vector and $\overline Z_0=\rho J_z\mathbf u$ are horizontal.
The product rule and $(P_T)_{hh}=I-B$ therefore give exactly
\[
 \begin{split}
 P_T:D\bigl((\phi\circ\pi)\overline Z_0\bigr)
 &=\phi P_T:D\overline Z_0
   -\frac{\rho^2}{r^2}\psi
       (K_z\mathbf u)\cdot(I-B)(J_z\mathbf u)\\
 &=\phi P_T:D\overline Z_0-\frac{\rho^2}{r^2}\psi
   +\frac{\rho^2}{r^2}\psi B[K_z\mathbf u,J_z\mathbf u].
 \end{split}
\]
The identity coefficient is one because $J_z,K_z$ are self-adjoint
inverses and $(K_z\mathbf u)\cdot(J_z\mathbf u)=|\mathbf u|^2=1$.
The radial terms are understood to be zero on the central fiber, as
their factors of $\rho^2$ require.

For the remaining derivative, the pointwise block calculation in the
proof of \Cref{lt:curvedvariation} gives
\[
 \begin{split}
 P_T:D\overline Z_0
 &=\Div_SZ_0-B:\nabla_SZ_0+\varepsilon_Z,\\
 |\varepsilon_Z|
 &\leq C\left\{|z|e|\nabla_SZ_0|
  +\left[(|\nabla^\perp H_S|+|z|)|z|+\sqrt e\,|z|\right]|Z_0|\right\}.
 \end{split}
\]
This is the estimate for the unweighted lift, so the product formula
multiplies its entire remainder by $\phi$.  Subtract the reference
identity $Q\int_S\Div_S(\phi Z_0)\,\dd\vol_S=0$, in which
$\Div_S(\phi Z_0)=\phi\Div_SZ_0-\rho^2\psi/r^2$.  Since
$\dot A=(2-m)A+\int^{(r)}\rho^2\psi\,\dd M$, we obtain, with
\[
 \widetilde R=\int^{(r)}\rho^2\psi
                  B[K_z\mathbf u,J_z\mathbf u]\,\dd\mu,
\]
the exact identity
\begin{equation}\label{in:horizontalidentity}
 \begin{split}
 \dot A={}&2A-E+\widetilde R
   +r^2\int^{(r)}\phi(\Div_SZ_0-m)\,\dd M\\
 &-r^2\int^{(r)}\phi B:(\nabla_SZ_0-I)\,\dd\mu
   +r^2\int^{(r)}\phi\varepsilon_Z\,\dd\mu.
 \end{split}
\end{equation}

In geodesic normal coordinates $y$, Gauss' lemma gives
$Z_0=y^i\partial_i$.  If $g^S_{ij}$ are the metric components, then
\[
 f:=\Div_SZ_0-m=y\cdot D\log\sqrt{\det(g^S_{ij})}.
\]
The logarithmic volume factor has vanishing first derivative at $0$
and bounded second derivative in the fixed chart.  Consequently
\[
 |f|\leq C\rho^2,\qquad |\nabla_Sf|\leq C\rho,
 \qquad \|a_r\|_{1,r}\leq K_{\mathrm{geo}},\quad a_r=f/r^2,
\]
for one fixed $K_{\mathrm{geo}}$.  Also
$\nabla_SZ_0=I+O(\rho^2)$.  The number
$K_{\mathrm{app}}=\max\{1,K_{\mathrm{geo}}\}$ bounds both actual
collar multipliers: $a_r$ here and $a=1$ in the bootstrap.  Choose
$C_\eta$ in the convention above to dominate $K_{\mathrm{app}}C_1$
and the finitely many other error prefactors in the frequency estimates.
Then $\eta\geq K_{\mathrm{app}}\eta_1$ has a fixed prefactor, so
\eqref{in:exponentialerror} applies.  The first line of
\eqref{in:collarestimates} and positivity of $B$ now give
\[
 \begin{split}
 r^2\left|\int^{(r)}\phi(\Div_SZ_0-m)\,\dd M\right|
   &\leq Cr^2(E+r^2H)+r^2\eta R,\\
 r^2\left|\int^{(r)}\phi B:(\nabla_SZ_0-I)\,\dd\mu\right|
   &\leq Cr^2E.
 \end{split}
\]
For the pointwise remainder use $|Z_0|\leq r$,
$|\nabla_SZ_0|\leq C$, and $\sup|z|\leq\delta r$.  The first
and last lines of \eqref{in:weightcomparison}, followed by
Cauchy--Schwarz, bound its three terms respectively by
\[
 \begin{split}
 r^2\int^{(r)}\phi|\varepsilon_Z|\,\dd\mu
 &\leq C\delta rE+Cr^3(H+\delta^2E)
                  +Cr^2\sqrt{E(H+\delta^2E)}\\
 &\leq Cr^2(E+H)+\eta E.
 \end{split}
\]
Here the term involving $\nabla^\perp H_S$ uses the weighted
curvature-square bound inside $\Lambda$; the last term comes from
$\sqrt e\,|z|\,|Z_0|$.

It remains to compare the radial stress $\widetilde R$ with $R$.
Positivity of $B$, its trace $e$, and $|J_z-K_z|\leq C|z|$ imply
$|B[v,(J_z-K_z)\mathbf u]|\leq C|z|\sqrt{be}$.  Thus
$|\widetilde R-R|$ is at most
\begin{equation}\label{in:radialreplacement}
 Cr\sqrt R\left(\int^{(r)}\psi|z|^2e\,\dd\mu\right)^{1/2},
\end{equation}
by Cauchy--Schwarz.  Since $e\leq m$, the second line of
\eqref{in:weightcomparison} gives
\[
 |\widetilde R-R|
 \leq Cr\sqrt{RH}+C\delta r\sqrt{R(E+R)}
 \leq Cr\sqrt{RH}+\eta(E+R).
\]
Combining these estimates in \eqref{in:horizontalidentity} and dividing
by $H$ proves \eqref{in:horizontalfrequency}.  The differentiated
kernel has been exhausted by $\widetilde R$; every other remainder in
that identity carries the weight $\phi$.

Finally, set
\[
 Y=\frac{\rho K_z\mathbf u}{|K_z\mathbf u|^2}.
\]
This field is defined independently of the varifold tangent plane.
Since $\nabla\rho=K_z\mathbf u$, it satisfies
$Y\cdot\nabla_T\rho=\rho h_r$.  Testing stationarity with
$\psi|z|^2Y$ consequently cancels the differentiated kernel exactly
and gives
\begin{equation}\label{in:heightidentity}
 \dot H=\int^{(r)}\psi\left\{
 |z|^2(\tr_T DY-mh_r)
 +\frac{2\rho z\cdot P_{vh}K_z\mathbf u}{|K_z\mathbf u|^2}
                         \right\}\,\dd\mu.
\end{equation}
There is no derivative of the measurable tangent projection in this
calculation.  Also $\nabla|z|^2=2z$ in the ambient Euclidean space,
so normal-frame connection terms do not arise from this factor.

The center estimates give $t(y)=o(\rho(y))$ and hence
$|z|\leq C\rho^2$ on the support near the origin.  Although
derivatives of $\mathbf u$ have size $O(\rho^{-1})$, differentiation
of $Y$ thus gives, on the support,
\[
 DY=P_h+O(\rho+|z|/\rho)=P_h+O(r).
\]
It follows that $\tr_TDY-mh_r=m(1-h_r)-e+O(r)$.
The positive $m(1-h_r)$ contribution is bounded by
$C\delta^2(E+R)$, by the proof of \eqref{in:weightcomparison}.
The $O(r)$ contribution is at most $CrH+\eta(E+R)$, and the term
$-e|z|^2$ can be discarded for an upper bound.  Replacing the
denominator in the second term of \eqref{in:heightidentity} by one
costs at most $C\delta r\sqrt{RH}\leq CrH+\eta R$.
Equation \eqref{in:normalflux} now proves the upper bound in
\eqref{in:heightfrequency}.

To justify the test at the central fiber, first remove $\rho<\epsilon$
by a smooth radial cutoff $\chi_\epsilon$, equal to zero for
$\rho\leq\epsilon$, equal to one for $\rho\geq2\epsilon$, and satisfying
$|\nabla\chi_\epsilon|\leq C/\epsilon$.  Since $|Y|\leq C\rho$,
$|Y||\nabla\chi_\epsilon|\leq C$ on the transition annulus.  Its extra
term is therefore bounded, at fixed $r$, by a constant times
$r^{-m}\int_{\rho<2\epsilon}|z|^2\,\dd\mu$, which tends to zero.
The preceding support bound and local mass bounds also justify all
other limiting terms.  The lower bound in
\eqref{in:heightfrequency} follows directly from differentiating
$r^{-m}\psi_r$: its derivative in $\log r$ is at least
$-m r^{-m}\psi_r$, and the remaining factor $|z|^2h_r$ is
nonnegative and independent of $r$.
\end{proof}

\subsection{Doubling and vanishing height}

\begin{proposition}[The inward bootstrap]\label{in:bootstrap}
On all sufficiently small radii, $P$ is bounded above and below.
For each fixed $\lambda\in(0,1)$ there are constants $C_\lambda,c_\lambda>0$
such that
\begin{equation}\label{in:doubling}
 c_\lambda H(r)\leq H(\lambda r)\leq C_\lambda H(r),
 \qquad
 \int_{\lambda r}^{r}q_*(v)\,\frac{\dd v}{v}\leq C_\lambda.
\end{equation}
There is a sequence $r_j\downarrow0$ on which both $n_*(r_j)$ and
$q_*(r_j)$ are bounded.
\end{proposition}

\begin{proof}
All quantities are locally absolutely continuous as functions of
$\log r$ where differentiated; $H>0$, and inequalities for derivatives
may be read almost everywhere.  Put
$v_0=\dot A/H$ and $h_0=\dot H/H$ temporarily.  The first line of
\eqref{in:collarestimates} with $a=1$, followed by
\eqref{in:normalfrequency}, gives
\[
 |P|\leq C(n_*+r^2)+\eta q_*
       \leq C(\sqrt{q_*}+r^2)+C\eta q_*.
\]
It follows, for a sufficiently large fixed $K$ and $|P|\geq K$, that
\begin{equation}\label{in:largefrequency}
 \frac{q_*}{|P|}\geq c\min\{|P|,\eta^{-1}\},
 \qquad v_0\geq cq_*.
\end{equation}
For the first assertion, either $C\sqrt{q_*}$ or $C\eta q_*$ is at least
a fixed fraction of $|P|$.  The second follows from
\eqref{in:horizontalfrequency}, increasing $K$ and then shrinking the
radius so that $q_*$ absorbs $2|P|$, $n_*$, and every error.

If $P\leq-K$, the lower bound $h_0\geq-m$ gives
\[
 \dot P=v_0-Ph_0\geq cq_*+mP>0.
\]
Such an excursion persists as the radius decreases: in inward time
$s=-\log r$, $P_s<0$ throughout this region.  There also $A<0$ and
$\dot A>0$, so $A$ decreases inward and stays below a fixed negative
number.  This contradicts $A\to0$, which follows from its factor
$r^2$ and bounded normalized mass.  We conclude that $P\geq-K$ on
the small-radius interval.

There are arbitrarily small radii at which $P\leq T^2$.  Otherwise,
eventually $P>T^2$, and \eqref{in:largefrequency} and
\eqref{in:exponentialerror} imply $v_0/P\geq cT^2$.  Define
\[
 Z=\log(r^2/A)=T-\log P.
\]
Inward differentiation gives $Z_s=v_0/P-2\geq cT^2-2$.
Since $T\geq2s-C$, $Z$ becomes positive and unbounded.  Also $Z\leq T$
in this regime, whence eventually $Z_s\geq c'Z^2$.
Integrating $(1/Z)_s\leq-c'$ forces a blow-up at finite inward time,
although $A,H$ are defined and positive throughout every finite
subinterval of this regime.  This is impossible.

We next prove the differential estimate
\begin{equation}\label{in:moderatefrequency}
 \dot P\geq-C(rP+gP^2)
                \qquad\hbox{when }K\leq P\leq T^4.
\end{equation}
Combining \eqref{in:horizontalfrequency} and
\eqref{in:heightfrequency} gives
\[
 \dot P\geq 2P-n_*+q_*-2Pn_*
 -C\{r\sqrt{q_*}+r^2(1+n_*)+rP\}
 -C\eta(1+P)(n_*+q_*).
\]
If $q_*\geq L P^2$ for a sufficiently large fixed $L$, then
$n_*\leq C\sqrt{q_*}+Cr^2$ and $\eta T^4=o(1)$ show that the right-hand
side is positive.  Otherwise $n_*\leq CP$, and squaring
\eqref{in:normalfrequency} gives
\[
 q_*\geq n_*^2-Cr^2P-C\eta P^2.
\]
The independent signed second line of \eqref{in:collarestimates}, with
$\eta_1\leq\eta$, gives
\[
 n_*\geq2P-d_0,\qquad
 d_0=C(gP+r^2+\eta P^2).
\]
For large $K$ and small $r$, this also gives $n_*>1$.  The principal
terms factor as
\[
 2P-n_*+n_*^2-2Pn_*=(n_*-1)(n_*-2P)\geq-CPd_0.
\]
Thus $\dot P\geq-C(rP+gP^2+\eta P^3)$.  The last term is absorbed
by $gP^2$, since $P\leq T^4$ and
$\eta T^4\leq g$ after shrinking the radius.  This proves
\eqref{in:moderatefrequency}.

It remains to justify that this restricted estimate gives a bound on
every smaller radius.  Choose an arbitrarily small moderate radius
$r_0=e^{-s_0}$, so $P_0\leq T_0^2$.  Take $s_0$ so large that
\[
 \int_{s_0}^{\infty}C\{e^{-s}+T(s)^{-16}\}\,\dd s
       \leq\varepsilon_0<1.
\]
This is possible by \eqref{in:exponentialerror}.  On each component
where $P\geq K$, before a possible exit through $P=T^4$,
\eqref{in:moderatefrequency} yields
\[
 (\log P)_s\leq C(r+gP)\leq C(r+T^{-16}).
\]
Moreover $Z_s>0$ there by \eqref{in:largefrequency}, increasing $K$
if necessary.  If the first component starts at $s_0$, these inequalities
give
\[
 P\leq e^{\varepsilon_0}P_0\leq e^{\varepsilon_0}T_0^2,
 \qquad
 T=Z+\log P\geq T_0-\log P_0+\log K\geq T_0/2.
\]
For large $T_0$ these bounds exclude $P=T^4$.  If a component starts
later, at $s_i$, then $P(s_i)=K$, and the same argument gives
$P\leq e^{\varepsilon_0}K$ and $T\geq T(s_i)$ there, again excluding
the upper exit.  This proves the finite uniform bound
\[
 -K\leq P(r)\leq e^{\varepsilon_0}\max\{K,P_0\}
                     \qquad(0<r\leq r_0).
\]

Write this bound as $|P|\leq B_0$.  After $B_0$ is fixed, choose
$0<\bar r_0\leq r_0$ so that $C\eta(1+B_0)\leq1/4$ for the constants in
\Cref{in:frequency} whenever $0<r\leq\bar r_0$; this absorbs its
$\eta q_*$ terms after quotient differentiation.  For $P\geq0$ use the
upper height derivative bound in $\dot P=v_0-Ph_0$; for $P<0$ use its
lower bound.  In both cases \Cref{in:frequency} and
$n_*\leq C\sqrt{q_*}+Cr^2$ imply
\begin{equation}\label{in:integralenergy}
 \dot P\geq cq_*-C_{B_0}.
\end{equation}
Integrating over a logarithmic interval of fixed length bounds
$\int q_*\,\dd\log r$.  Cauchy--Schwarz then bounds
$\int n_*\,\dd\log r$, and integration of
\eqref{in:heightfrequency} proves both sides of the height comparison
in \eqref{in:doubling}.  In each interval
$[2^{-j-1}\bar r_0,2^{-j}\bar r_0]$ choose a radius where $q_*$ is
bounded by its logarithmic average.  Equation \eqref{in:normalfrequency}
bounds $n_*$ at those radii, proving the last assertion.
\end{proof}

\begin{proposition}[Vanishing height at a typical point]\label{in:finish}
The positive-height alternative at the origin is impossible.
Consequently every point satisfying the typical-point hypotheses \eqref{ce:typical} is regular.
\end{proposition}

\begin{proof}
Take the radii $r_j$ from \Cref{in:bootstrap}, and write $r=r_j$.
On $\rho<3r/4$ the weights $\phi,\psi$ have positive lower bounds
depending only on $p$.  Equations \eqref{in:weightcomparison},
\eqref{in:normalfrequency}, and the boundedness of $n_*,q_*$ give
\begin{equation}\label{in:selectedenergy}
 r^{-m}\int_{\rho<3r/4}|z|^2\,\dd\mu
       +r^{2-m}(D+\Lambda)(B_{3r/4})\leq CH(r).
\end{equation}
Here a height integral is over the part of the separated tube above
the indicated base ball, whereas $D,\Lambda$ are projected measures.
On the other hand, applying \eqref{in:doubling} with $\lambda=1/4$
and using $\psi\leq2p$, $h_r\leq1$ at radius $r/4$ gives
\begin{equation}\label{in:heightlower}
 r^{-m}\int_{\rho<r/4}|z|^2\,\dd\mu\geq cH(r)>0.
\end{equation}

Fix an exceptional fraction $\delta_{\mathrm{cfg}}=1/4$ and a finite
exponent $p_*>2$ satisfying $m(1/2-1/p_*)<1$.  Use the enlargement $c$
fixed from the local data in \Cref{lt:config}, with the required
scalar-selection room, and then fix
$d_{\mathrm{cfg},0}=d_0(\delta_{\mathrm{cfg}})>0$.  Then choose
$\epsilon>0$ so small that $B_{r/4}$ has a finite cover by balls of
radius $\epsilon r$ whose required enlargements lie in $B_{r/2}$.
The number of balls and $\epsilon$ are fixed independently of $r$.

Apply \Cref{lt:scalar} with room between $3r/4$ and $r/2$; its
integer is again $Q$.  Its variation estimate and
\eqref{in:selectedenergy} give, for every enlarged ball
$B_{c\epsilon r}(x)$ of this cover,
\[
 (\epsilon r)^{-m}\tau(B_{c\epsilon r}(x))
       \leq C_\epsilon H(r)/r^2\longrightarrow0,
\]
uniformly across the cover.  The normalized smooth mass averages
used in the local scalar selections differ from $Q$ by at most the
same bound.  Thus, after discarding finitely many selected radii, the
scalar smallness and smooth-average conditions select $Q$ on every
required local ball, and its displayed budget is below
$d_{\mathrm{cfg},0}$.  Applying \Cref{lt:config} with the fixed
choices above now supplies on each covering ball a configuration with
the prescribed majority guarantee and a configuration-distance moment
whose normalized $L^{p_*}$ norm is at most $C_\epsilon\sqrt{H(r)}$.

Every entry of each such configuration has size at most
$C_\epsilon\sqrt{H(r)}$.  Indeed, on a fixed majority of its base
ball the entries match the actual heights within this error.  The
area formula, the bounded projection Jacobian, and
\eqref{in:selectedenergy} show that a subset of this majority has
sum of squared actual heights at most $C_\epsilon H(r)$.  Every
configuration entry has positive integer weight, proving the claim.
Adding these configuration heights to the distance moment and
summing the fixed finite cover yields
\begin{equation}\label{in:highermoment}
 r^{-m}\int_{\rho<r/4}|z|^{p_*}\,\dd\mu
                         \leq CH(r)^{p_*/2}.
\end{equation}

Let $\mathcal I=\{t>0\}$ on the center.  By \Cref{ce:center}, its
relative base volume tends to zero and the height vanishes above
$S\setminus\mathcal I$.  The same is true for normalized projected
mass, since \Cref{lt:scalar} gives
\begin{equation}\label{in:influencemass}
 \begin{split}
 r^{-m}\pi_\#\mu(\mathcal I\cap B_{r/4})
 &\leq Qr^{-m}\vol_S(\mathcal I\cap B_{r/4})
                  +r^{-m}|M|(B_{r/4})\\
 &\leq o(1)+CH(r)/r^2=o(1).
 \end{split}
\end{equation}
H\"older's inequality with \eqref{in:highermoment} and
\eqref{in:influencemass} now gives
\[
 r^{-m}\int_{\rho<r/4}|z|^2\,\dd\mu
 \leq CH(r)
   \left[r^{-m}\pi_\#\mu(\mathcal I\cap B_{r/4})\right]^{1-2/p_*}
 =o(H(r)),
\]
contradicting \eqref{in:heightlower}.  Thus the varifold lies on the
center near the origin, and \Cref{ce:constancy} gives regularity.
\end{proof}

\begin{proof}[Completion of the proof of \Cref{thm:main}]
The weighted rectifiable tangent and density theorem
\cite[Chapter~3, Definition~1.7, Remark~1.8, and Theorem~1.9]{Simon2018},
together with differentiation of the integer multiplicity sets,
imply that at $\mu$-almost every point $x$ there is a positive integer $Q$ and
an approximate tangent plane such that the dilated varifolds tend to
the multiplicity-$Q$ plane and
\[
 \frac{\mu(\{y:\theta(y)\ne Q\}\cap \mathbf B_r(x))}{r^m}
                         \longrightarrow0.
\]
The tangent statement includes convergence of the tangent projections;
it follows by differentiating the measurable tangent-plane map on a
countable rectifiable covering.  These are precisely the typical-point
hypotheses used in the center construction.  No regularity conclusion
is used in choosing the point.  \Cref{in:finish} therefore shows that
$\mu(\operatorname{Sing}(V))=0$.

For completeness, this implies the requested Hausdorff-measure
conclusion even for support points outside the chosen representative
of the rectifiable set.  Let $N\subset\spt\mu$ be a $\mu$-null set
in a compact subset of $U$, and let $O$ be an open neighborhood of
$N$ with arbitrarily small $\mu(O)$.  For each $x\in N$ take arbitrarily
small balls $\mathbf B_r(x)\subset O\cap U$.  The support-point
monotonicity bound in \Cref{lt:monotonicity} gives
$\mu(\mathbf B_r(x))\geq\omega_m r^m$.  The disjoint-ball covering
lemma selects pairwise disjoint such balls whose fivefold enlargements
cover $N$.  Their radii can all be taken below any prescribed bound,
and therefore
\[
 \HH^m_{\epsilon}(N)
       \leq C_m\sum_i r_i^m
       \leq C_m\sum_i\mu(\mathbf B_{r_i}(x_i))
       \leq C_m\mu(O).
\]
Let $\mu(O)\downarrow0$ and then $\epsilon\downarrow0$.
An exhaustion of $U$ proves $\HH^m\ll\mu$ on $\spt\mu$.
Applying this to the singular set proves
$\HH^m(\operatorname{Sing}(V))=0$.
\end{proof}
